\chapter{The heroes}
\ficha{heroe_barbaro}{The barbarian}{%
\begin{tabular}{@{}lc@{\qquad}lc@{}}
Body points &8& Attack dice &3\\ Mind points &2& Defence dice &2\\
\end{tabular}}{%
Default name: \textbf{Sigmar}. The party's most fearsome warrior: he withstands the most blows and strikes back hardest. He knows no magic, and his limited mental strength makes him vulnerable to spells, but he can wield any weapon and wear any armour. Put him at the front, opening doors and facing whatever waits behind them.}
\ficha{heroe_enano}{The dwarf}{%
\begin{tabular}{@{}lc@{\qquad}lc@{}}
Body points &7& Attack dice &2\\ Mind points &3& Defence dice &2\\
\end{tabular}}{%
Default name: \textbf{Grugni}. A good fighter, almost as tough as the barbarian, and a trap expert. He starts with a tool kit to disarm discovered traps. He cannot use magic either. A good weapon bought with his first gold makes him a formidable combatant.}
\ficha{heroe_duende}{The elf}{%
\begin{tabular}{@{}lc@{\qquad}lc@{}}
Body points &6& Attack dice &2\\ Mind points &4& Defence dice &2\\
\end{tabular}}{%
Default name: \textbf{Ladril}. The elf combines warrior and magician. He fights like the dwarf, can use all weapons and armour, and has one group of three spells chosen at the start of each adventure.}
\ficha{heroe_hechicero}{The wizard}{%
\begin{tabular}{@{}lc@{\qquad}lc@{}}
Body points &4& Attack dice &1\\ Mind points &6& Defence dice &2\\
\end{tabular}}{%
Default name: \textbf{Zoltan}. Weak in close combat and with very few body points, he makes up for his frailty with three spell groups---nine spells in all---and the strongest mind in the party. Only the staff and tool kit are useful shop purchases for him: the staff allows diagonal attacks from behind a companion. Protect him: spells used at the right moment decide adventures.}
\begin{figure}[H]\centering
\includegraphics[width=0.48\linewidth]{heroe_barbaro_anim}\hfill
\includegraphics[width=0.48\linewidth]{heroe_enano_anim}\\[3pt]
\includegraphics[width=0.48\linewidth]{heroe_duende_anim}\hfill
\includegraphics[width=0.48\linewidth]{heroe_hechicero_anim}
\caption{The four heroes in motion.}\end{figure}

\chapter{The enemies}
Morcar commands eight types of creature. The table summarises their characteristics; movement is the number of squares they can travel in a turn.
\begin{table}[H]\centering\small
\begin{tabular}{@{}llcccc@{}}
\toprule\textbf{Enemy} & \textbf{Screen name} & \textbf{Move} & \textbf{Attack} & \textbf{Defence} & \textbf{Body}\\\midrule
Goblin &\pant{GOBLIN}&10&2&1&1\\
Orc &\pant{ORC}&8&3&2&1\\
Fimir &\pant{FIMIR}&6&3&3&1\\
Chaos warrior &\pant{WARRIOR}&6&3&4&1\\
Skeleton &\pant{SKELETON}&6&2&2&1\\
Zombie &\pant{ZOMBIE}&4&2&3&1\\
Mummy &\pant{MUMMY}&4&3&4&1\\
Gargoyle &\pant{SORCERER}&6&4&4&1\\\bottomrule
\end{tabular}\caption{Enemy characteristics: attack and defence values are dice.}\end{table}
Remember that monsters defend on only one of the six die faces: even the best protected eventually fall.
\ficha{monstruo_goblin}{Goblin}{Movement 10 \quad Attack 2 \quad Defence 1}{%
Small, cowardly and very fast. Their screen name is \pant{GOBLIN}. One alone is little threat, but they arrive first and usually in groups.}
\ficha{monstruo_orco}{Orc}{Movement 8 \quad Attack 3 \quad Defence 2}{%
Morcar's infantry: strong, numerous and always ready to fight. They hit as hard as the barbarian. The heroes will meet several of their chieftains.}
\ficha{monstruo_fimir}{Fimir}{Movement 6 \quad Attack 3 \quad Defence 3}{%
Huge one-eyed reptiles armed with axes and maces. Slower than orcs, but considerably harder to bring down.}
\ficha{monstruo_guerrero_caos}{Chaos warrior}{Movement 6 \quad Attack 3 \quad Defence 4}{%
Morcar's champions, clad in heavy black armour. They have the best defence among the living, and often command the other creatures and guard important locations.}
\ficha{monstruo_esqueleto}{Skeleton}{Movement 6 \quad Attack 2 \quad Defence 2}{%
The most common undead, raised from their graves by Chaos magic. Fragile but tireless. Like all undead, they fear holy water and the Spirit Blade.}
\ficha{monstruo_zombi}{Zombie}{Movement 4 \quad Attack 2 \quad Defence 3}{%
Animated corpses, slow and clumsy. They are easy to leave behind, but withstand more than their appearance suggests.}
\ficha{monstruo_momia}{Mummy}{Movement 4 \quad Attack 3 \quad Defence 4}{%
Guardians of ancient tombs. As slow as zombies, but they strike like orcs and withstand blows like Chaos warriors.}
\ficha{monstruo_gargola}{Gargoyle}{Movement 6 \quad Attack 4 \quad Defence 4}{%
Morcar's most powerful creature: a stone demon with the game's best attack and defence. The same figure represents Chaos sorcerers, hence its screen name \pant{SORCERER}.}
\subsection{Named enemies}
Some adventures have a main opponent who is more dangerous than ordinary creatures, often with a reward on their head:
\begin{description}
\item[Ulag] Orc warlord responsible for Sir Ragnar's abduction.
\item[Grak] Ulag's son, sworn to avenge his father.
\item[Gulthor] Chaos warrior leading the band that stole Prince Magnus' gold.
\item[Balur] The Fire Mage. No fire spell can hurt him.
\item[The Witch Lord] Shown as \pant{WITCHLORD}, the King of the Dead. Immune to spells; of all weapons, only the Spirit Blade can wound him.
\end{description}
\textit{Return of the Witch Lord} also features Bellthor, the guardian gargoyle; Skulmar, captain of the Forgotten Legion; Kessandria, the Witch Queen; the Spirit Riders; and the Doomguard warriors, all stronger than ordinary creatures.
\subsection{Dungeon furniture}
Rooms are not empty: tables, torture racks, alchemists' benches, tombs, fireplaces, shelves, weapon racks, chests and thrones occupy squares and block passage. You cannot walk over them, but they often identify rooms worth searching.
\begin{figure}[H]\centering
\includegraphics[width=\linewidth]{mobiliario_a}\\[4pt]
\includegraphics[width=\linewidth]{mobiliario_b}\\[6pt]
\includegraphics[width=0.12\linewidth]{monstruo_cofre}\qquad
\includegraphics[width=0.12\linewidth]{monstruo_trono}
\caption{Furniture, chest and throne.}\end{figure}
